\documentclass[a4paper,9pt]{article}

% --------------------------------------------------
% Packages
% --------------------------------------------------
\usepackage[margin=0.9cm]{geometry}
\usepackage[T1]{fontenc}
\usepackage[utf8]{inputenc}
\usepackage{lmodern}
\usepackage{amsmath,amssymb}
\usepackage{xcolor}
\usepackage{multicol}
\usepackage{enumitem}
\usepackage[most]{tcolorbox}
\usepackage{titlesec}
\usepackage{microtype}
\usepackage{booktabs}

% --------------------------------------------------
% Colors
% --------------------------------------------------
\definecolor{unibaselteal}{RGB}{165,215,210}
\definecolor{darkteal}{RGB}{20,90,85}
\definecolor{softgray}{RGB}{245,245,245}
\definecolor{darkgray}{RGB}{70,70,70}
\definecolor{rmsred}{RGB}{220,0,0}
\definecolor{adaorange}{RGB}{210,130,0}
\definecolor{adampurple}{RGB}{130,70,170}

% --------------------------------------------------
% Formatting
% --------------------------------------------------
\setlength{\parindent}{0pt}
\setlength{\parskip}{2pt}
\setlength{\columnsep}{0.55cm}

\setlist[itemize]{leftmargin=1.15em,itemsep=1pt,topsep=1pt}
\setlist[enumerate]{leftmargin=1.3em,itemsep=1pt,topsep=1pt}

\titleformat{\section}
  {\large\bfseries\color{darkteal}}
  {}
  {0pt}
  {}

\titleformat{\subsection}
  {\normalsize\bfseries}
  {}
  {0pt}
  {}

\tcbset{
  sharp corners,
  boxrule=0.4pt,
  colback=softgray,
  colframe=unibaselteal,
  left=4pt,
  right=4pt,
  top=3pt,
  bottom=3pt
}

\newtcolorbox{qa}[1]{
  enhanced,
  breakable,
  colback=white,
  colframe=unibaselteal,
  title=\textbf{#1},
  fonttitle=\small\bfseries,
  coltitle=black,
  attach boxed title to top left={xshift=2mm,yshift=-2mm},
  boxed title style={colback=unibaselteal,colframe=unibaselteal},
  before skip=5pt,
  after skip=5pt,
  left=4pt,
  right=4pt,
  top=5pt,
  bottom=3pt
}

\newcommand{\key}[1]{\textbf{\color{darkteal}#1}}
\newcommand{\method}[1]{\textbf{#1}}
\newcommand{\oral}{\textit{\color{darkgray}Oral answer: }}

% --------------------------------------------------
% Document
% --------------------------------------------------
\begin{document}

\begin{center}
    {\LARGE \textbf{Q\&A Cheatsheet --- Adaptive Algorithms}}\\
    \vspace{2pt}
    {\large Continuous Optimization Project Presentation}\\
    \vspace{2pt}
    Erion Krasniqi \& Denis Mustafa Xhabrahimi
\end{center}

\vspace{-2mm}

\begin{tcolorbox}[colback=unibaselteal!35,colframe=unibaselteal]
\textbf{Golden Rule for oral Q\&A:}
Answer in \textbf{3 sentences}: 
\textbf{intuition} $\rightarrow$ \textbf{project connection} $\rightarrow$ \textbf{result / consequence}.\\
Avoid long formulas unless explicitly asked. Start intuitive, then become precise if needed.
\end{tcolorbox}

\begin{multicols}{2}

% ==================================================
\section{1. Big Picture}

\begin{qa}{Q: What is the main goal of your project?}
\oral We compare adaptive first-order optimization methods against SGD. 
We focus on Adagrad, AdaGrad-Norm, RMSprop and Adam. 
Experimentally, we test whether they improve convergence speed and final performance on A9a and MNIST.
\end{qa}

\begin{qa}{Q: What is your main takeaway?}
\oral Adaptive methods converge much faster in early iterations. 
However, final test accuracy is often very similar across optimizers. 
So the main benefit is \textbf{optimization speed}, not necessarily better final generalization.
\end{qa}

\begin{qa}{Q: Why adaptive step sizes?}
\oral A fixed learning rate is difficult because one value may be too large for some directions and too small for others. 
Adaptive methods use gradient history to rescale the step size automatically. 
This makes training less sensitive and often much faster at the beginning.
\end{qa}

\begin{qa}{Q: Why is SGD your baseline?}
\oral SGD is the standard stochastic first-order method. 
It is simple, cheap per iteration and well understood theoretically. 
Therefore it is the natural baseline for checking whether adaptive methods really help.
\end{qa}

% ==================================================
\section{2. Theory Basics}

\begin{qa}{Q: What is the basic gradient descent update?}
\oral Gradient descent updates the parameter vector in the negative gradient direction. 
The gradient gives the direction of steepest increase, so we move in the opposite direction. 
The learning rate controls how large this movement is.
\[
x_{t+1}=x_t-\eta\nabla f(x_t)
\]
Here, \(x_t\) is the current point, \(\nabla f(x_t)\) is the gradient, and \(\eta\) is the step size.
\end{qa}

\begin{qa}{Q: What is the theoretical step-size choice?}
\oral For an \(L\)-smooth function, theory often requires a step size such as \(\eta\leq 1/L\). 
This prevents the update from being too aggressive. 
In practice, however, \(L\) is often unknown or too conservative.
\[
\eta \leq \frac{1}{L}
\]
Here, \(L\) measures how quickly the gradient can change.
\end{qa}

\begin{qa}{Q: What does convexity mean?}
\oral Convexity means the loss landscape has no bad local minima. 
If we reach a stationary point in a convex problem, we have reached a global optimum. 
This makes convex logistic regression a clean baseline for comparing optimizers.
\[
f(\lambda x+(1-\lambda)y)
\leq
\lambda f(x)+(1-\lambda)f(y)
\]
\end{qa}

\begin{qa}{Q: Why is logistic loss convex?}
\oral Logistic loss has non-negative curvature everywhere. 
Its Hessian can be written as \(X^\top D X\), where \(D\) has non-negative diagonal entries. 
Therefore the Hessian is positive semidefinite, so the loss is convex.
\[
H=X^\top D X
\]
For any vector \(z\):
\[
z^\top H z = z^\top X^\top D Xz = (Xz)^\top D(Xz) \geq 0
\]
because \(D\) has only non-negative diagonal entries.
\end{qa}

\begin{qa}{Q: What does positive semidefinite mean?}
\oral A matrix is positive semidefinite if it has no negative curvature. 
Formally, \(z^\top H z\geq 0\) for all vectors \(z\). 
For optimization, this means the function is convex if the Hessian is positive semidefinite everywhere.
\[
H\succeq 0
\quad \Longleftrightarrow \quad
z^\top H z\geq 0
\]
\end{qa}

% ==================================================
\section{3. Loss Functions \& Formulas}

\begin{qa}{Q: What is the logistic loss formula?}
\oral Logistic loss is the negative log-likelihood for binary classification.
It is convex and smooth --- ideal for testing convergence guarantees.

\[
f(w) = \frac{1}{n}\sum_{i=1}^n \log\bigl(1 + \exp(-y_i \cdot w^\top x_i)\bigr)
\]

\textbf{Gradient:}
\[
\nabla f(w) = -\frac{1}{n}\sum_{i=1}^n \frac{y_i\, x_i}{1 + \exp(y_i\, w^\top x_i)}
\]

\textbf{Hessian} (proves convexity):
\[
\nabla^2 f(w) = \frac{1}{n} X^\top D X, \quad D_{ii} = \sigma_i(1-\sigma_i) > 0
\]
where \(\sigma_i = \sigma(y_i w^\top x_i)\) is the sigmoid.
\end{qa}

\begin{qa}{Q: What is softmax cross-entropy?}
\oral Multi-class generalization of logistic loss.
Measures how well predicted probabilities match true labels.
Convex in the weight matrix \(W\).

\[
f(W) = -\frac{1}{n}\sum_{i=1}^n \log \frac{\exp(w_{y_i}^\top x_i)}{\sum_{k=1}^K \exp(w_k^\top x_i)}
\]

where \(W = [w_1, \ldots, w_K]\), \(K=10\) classes.

\textbf{Gradient} w.r.t.\ class \(k\):
\[
\nabla_{w_k} f = \frac{1}{n}\sum_{i=1}^n (p_{ik} - \mathbf{1}[y_i=k])\, x_i
\]

with \(p_{ik} = \text{softmax}_k(Wx_i) = \frac{\exp(w_k^\top x_i)}{\sum_j \exp(w_j^\top x_i)}\).
\end{qa}

% ==================================================
\section{4. Non-Convex Regularization}

\begin{qa}{Q: What makes your regularizer non-convex?}
\oral For small weights, the regularizer behaves similarly to an \(\ell_2\) penalty. 
For larger weights, it saturates instead of growing indefinitely. 
This curvature change introduces non-convexity.
\[
R(w)
=
\lambda
\sum_j
\frac{\alpha w_j^2}
{1+\alpha w_j^2}
\]
\end{qa}

\begin{qa}{Q: What is the intuition behind this regularizer?}
\oral The term penalizes small and medium weights, but the penalty becomes bounded for large weights. 
So very large weights are not punished infinitely strongly. 
This creates a softer regularization effect compared to standard \(\ell_2\).
\[
\frac{\alpha w_j^2}{1+\alpha w_j^2}
\rightarrow 1
\quad \text{for large } |w_j|
\]
\end{qa}

\begin{qa}{Q: Why does the non-convex regularizer have limited impact?}
\oral At moderate \(\lambda\), the regularization term is small compared to logistic loss. 
So the overall landscape is only mildly non-convex. 
That is why the optimizer ranking stays almost the same.
\[
F(w)= \text{logistic loss} + \lambda R(w)
\]
If \(\lambda\) is small, the second term only mildly changes the landscape.
\end{qa}

\begin{qa}{Q: What happens if \(\lambda\) becomes larger?}
\oral If \(\lambda\) becomes larger, the regularizer starts to dominate. 
Then it pushes weights toward smaller values and can cause underfitting. 
In our \(\lambda\)-sweep, accuracy drops when \(\lambda \geq 0.1\).
\end{qa}

\begin{qa}{Q: How is \(\lambda\) chosen?}
\oral \(\lambda\) is a hyperparameter. 
It controls the trade-off between data fit and regularization. 
We use \(\lambda = 0.01\), \(\alpha = 1\) following Reddi et al.\ (2016).
\end{qa}

% ==================================================
\section{5. Parameter Settings}

\begin{tcolorbox}[title=\textbf{All Hyperparameters Used}, colback=white]
\renewcommand{\arraystretch}{1.15}
\begin{tabular}{@{}l l l@{}}
\toprule
\textbf{Parameter} & \textbf{Value} & \textbf{Note} \\
\midrule
\multicolumn{3}{@{}l}{\textit{Shared:}} \\
Batch size (A9a) & 128 & \\
Batch size (MNIST) & 256 & \\
Epochs & 10--15 & \\
Seeds & 3 & averaged \\
\(\epsilon\) & \(10^{-8}\) & all adaptive \\
\midrule
\multicolumn{3}{@{}l}{\textit{SGD:}} \\
\(\eta\) & tuned & grid search \\
\midrule
\multicolumn{3}{@{}l}{\textit{Adagrad (Duchi et al.\ 2011):}} \\
\(\eta\) & tuned & base LR \\
\(G_0\) & 0 & per-coord accum. \\
\midrule
\multicolumn{3}{@{}l}{\textit{AdaGrad-Norm (Ward et al.\ 2020):}} \\
\(\eta\) & tuned & base LR \\
\(b_0^2\) & 0 & scalar accum. \\
\midrule
\multicolumn{3}{@{}l}{\textit{RMSprop (Tieleman \& Hinton 2012):}} \\
\(\eta\) & tuned & base LR \\
\(\rho\) & 0.9 & decay factor \\
\midrule
\multicolumn{3}{@{}l}{\textit{Adam (Kingma \& Ba 2014):}} \\
\(\eta\) & tuned & base LR \\
\(\beta_1\) & 0.9 & 1st moment decay \\
\(\beta_2\) & 0.999 & 2nd moment decay \\
\midrule
\multicolumn{3}{@{}l}{\textit{Non-convex regularizer:}} \\
\(\lambda\) & 0.01 & reg.\ strength \\
\(\alpha\) & 1 & curvature param \\
\bottomrule
\end{tabular}
\end{tcolorbox}

\begin{qa}{Q: How did you choose learning rates?}
\oral Log-scale sweep: \(\eta \in \{10^{-4}, 3\!\times\!10^{-4}, 10^{-3}, \ldots, 1\}\). 
Picked best final training loss per optimizer.
Key: Adagrad very robust across range; RMSprop diverges at \(\eta=1\).
\end{qa}

\begin{qa}{Q: Why \(\beta_2 = 0.999\) for Adam?}
\oral Default from Kingma \& Ba (2014). Very conservative --- needs \(\sim\)1000 iterations to stabilize \(\hat{v}_t\). 
This is why Adam adapts more slowly initially vs.\ RMSprop (\(\rho=0.9\), adapts in \(\sim\)10 steps).
\end{qa}

\begin{qa}{Q: Why \(\rho = 0.9\) for RMSprop?}
\oral Standard from Hinton's lecture notes. 
Effective memory \(\sim\)10 steps (\(0.9^{10}\approx 0.35\)). 
Much faster than Adam's \(\beta_2=0.999\) (memory \(\sim\)1000 steps).
\end{qa}

\begin{qa}{Q: Why mini-batches of 128/256?}
\oral Standard compromise: full-batch is stable but slow, single-sample is fast but noisy. 
128 for A9a (smaller), 256 for MNIST (larger). 
Matches stochastic optimization setting of our convergence guarantees.
\end{qa}

% ==================================================
\section{6. Convergence Rates}

\begin{qa}{Q: What does \(O(1/\sqrt{T})\) mean intuitively?}
\oral Sublinear convergence. Progress slows over time. 
To halve the error: need \(4\times\) more iterations.
\[
\text{error}(T)\sim \frac{1}{\sqrt{T}}
\qquad
T_{\text{new}}=4T \text{ to halve error}
\]
\end{qa}

\begin{qa}{Q: Is \(O(1/\sqrt{T})\) linear convergence?}
\oral No. Linear = error multiplied by constant factor each step (\(e_t \leq c^t\), \(0<c<1\)). 
\(O(1/\sqrt{T})\) is sublinear --- flattens out.
On log-plot: linear = straight line; sublinear = curves.
\end{qa}

\begin{qa}{Q: Difference between convergence rate and guarantee?}
\oral Rate = how fast (e.g.\ \(O(1/\sqrt{T})\)). 
Guarantee = under which assumptions (smoothness, unbiased gradients, bounded variance). 
RMSprop: no guarantee; SGD: \(O(1/\sqrt{T})\) under standard assumptions.
\end{qa}

\begin{qa}{Q: Why are empirical rates faster than theoretical?}
\oral Theory = worst-case for arbitrary/adversarial functions. 
Our logistic/softmax have nice structure (smooth, bounded spectral norm). 
So empirical convergence is much faster than the bounds predict.
\end{qa}

% ==================================================
\section{7. Polyak Step Size}

\begin{qa}{Q: What is the Polyak step size?}
\oral Adapts the step size based on the function value gap to the optimum.
Derived from minimizing a quadratic upper bound along \(-\nabla f(x_t)\).
Requires knowledge of \(f^*\) (often unknown in practice).

\[
\eta_t = \frac{f(x_t) - f^*}{\|g_t\|^2}, \quad g_t = \nabla f(x_t)
\]

\textbf{Derivation:} For \(L\)-smooth convex \(f\):
\[
f(x) \leq f(x_t) + \nabla f(x_t)^\top(x-x_t) + \frac{L}{2}\|x-x_t\|^2
\]
Setting this model \(= f^*\) and solving along \(-\nabla f(x_t)\) gives \(\eta_t\).

\textbf{Remark:} The factor 2 from Taylor is absorbed because Polyak minimizes the quadratic model along the gradient direction (not arbitrary directions).
\end{qa}

\begin{qa}{Q: Why does Polyak adapt automatically?}
\oral \textbf{Near \(x^*\):} gap small \(\Rightarrow\) small steps (no overshooting).\\
\textbf{Far from \(x^*\):} gap large, but \(\|g_t\|\) also large \(\Rightarrow\) self-regulates.\\
\textbf{Far + flat:} \(\|g_t\|\) small \(\Rightarrow\) large step (needs big jump on flat terrain).

Adapts to distance-to-optimum \textbf{and} local gradient magnitude simultaneously.
\end{qa}

\begin{qa}{Q: Polyak algorithm?}
\oral
\begin{enumerate}[nosep]
\item Compute \(g_t = \nabla f(x_t)\)
\item Set \(\eta_t = (f(x_t) - f^*) / \|g_t\|^2\)
\item Update \(x_{t+1} = x_t - \eta_t g_t\)
\end{enumerate}
Requires: smooth convex \(f\), known optimal value \(f^*\).
\end{qa}

\begin{qa}{Q: Why didn't you include Polyak?}
\oral (1) Requires \(f^*\) --- unknown for our problems.
(2) Different adaptivity family: function-value-based vs.\ gradient-accumulation.
(3) Our scope: methods using \textbf{only gradient history}.
\end{qa}

\begin{qa}{Q: Connection to your project?}
\oral Both Polyak and Adagrad/Adam adapt the step size. 
Polyak adapts via \emph{distance to optimal} (needs \(f^*\)). 
Adagrad/Adam adapt via \emph{gradient behavior} (needs only gradient samples --- more practical when \(f^*\) unknown).
\end{qa}

% ==================================================
\section{8. SGD vs Adaptive Methods}

\begin{qa}{Q: Difference between SGD and adaptive methods?}
\oral SGD uses one global \(\eta\) for all coordinates. 
Adaptive methods change effective step size using past gradient info. 
Helps when coordinates have different scales.
\[
w_{t+1} = w_t-\eta g_t \quad \text{(same \(\eta\) everywhere)}
\]
\end{qa}

\begin{qa}{Q: Typical assumptions for SGD's \(O(1/\sqrt{T})\)?}
\oral (1) \(L\)-smooth: \(\|\nabla f(x)-\nabla f(y)\| \leq L\|x-y\|\).
(2) Unbiased: \(\mathbb{E}[g_t]=\nabla f(w_t)\).
(3) Bounded variance: \(\mathbb{E}\|g_t-\nabla f(w_t)\|^2 \leq \sigma^2\).
\end{qa}

% ==================================================
\section{9. Adagrad}

\begin{qa}{Q: How does Adagrad work?}
\oral Scales update using accumulated past squared gradients (per coordinate).
Large past gradients \(\Rightarrow\) smaller effective step size.

\[
G_t = G_{t-1} + g_t \odot g_t
\qquad
w_{t+1} = w_t - \frac{\eta}{\sqrt{G_t}+\epsilon}\, g_t
\]

Effective LR: \(\eta_t^{\text{eff}} = \eta/(\sqrt{G_{t,j}}+\epsilon)\). Shrinks monotonically.
Rate: \(O(\log T/\sqrt{T})\).
\end{qa}

\begin{qa}{Q: Intuition behind Adagrad?}
\oral More accumulated gradient magnitude \(\Rightarrow\) more cautious steps.
\[
G_t \uparrow \;\Rightarrow\; \sqrt{G_t} \uparrow \;\Rightarrow\; \eta_t^{\text{eff}} \downarrow
\]
Stabilizes optimization; avoids too aggressive steps on already-explored directions.
\end{qa}

\begin{qa}{Q: Why can Adagrad slow down?}
\oral \(G_t\) only grows (\(g_t^2 \geq 0\), so \(G_t \geq G_{t-1}\)). 
Over time: \(\eta/\sqrt{G_t} \to 0\). 
Good for convex (need decreasing steps), bad for long non-convex training.
\end{qa}

% ==================================================
\section{10. AdaGrad-Norm}

\begin{qa}{Q: How does AdaGrad-Norm work?}
\oral One global scalar accumulator (not per-coordinate).
Accumulates squared norm of full gradient.

\[
b_t^2 = b_{t-1}^2 + \|g_t\|^2
\qquad
w_{t+1} = w_t - \frac{\eta}{b_t+\epsilon}\, g_t
\]

Same \(O(1/\sqrt{T})\) rate. Cleaner non-convex theory (Ward et al.\ 2020).
\end{qa}

\begin{qa}{Q: Difference Adagrad vs.\ AdaGrad-Norm?}
\oral Adagrad: per-coordinate \(G_{jj}\) --- each feature gets own LR. Good for heterogeneous/sparse features.\\
AdaGrad-Norm: one scalar \(b_t\) for all --- simpler, loses per-feature adaptation.\\
MNIST (784 heterogeneous features): Adagrad wins clearly.
\end{qa}

% ==================================================
\section{11. RMSprop}

\begin{qa}{Q: How does RMSprop work?}
\oral Exponential forgetting of old squared gradients (moving average instead of sum).

\[
v_t = \rho\, v_{t-1} + (1-\rho)\,g_t^2
\qquad
w_{t+1} = w_t - \frac{\eta}{\sqrt{v_t}+\epsilon}\, g_t
\]

\(\rho=0.9\): keeps 90\% old, adds 10\% new.
\(v_t\) can increase \textbf{or decrease} (unlike Adagrad's \(G_t\)).
No formal convergence guarantee.
\end{qa}

\begin{qa}{Q: Why is RMSprop different from Adagrad?}
\oral Adagrad: \(G_t = G_{t-1} + g_t^2\) --- only grows, never forgets.\\
RMSprop: \(v_t = \rho v_{t-1} + (1-\rho)g_t^2\) --- forgets old info.\\
Avoids permanent LR decay \(\Rightarrow\) keeps adapting during long training.
\end{qa}

\begin{qa}{Q: Why can RMSprop diverge at large \(\eta\)?}
\oral If \(v_t\) becomes small (gradients recently small): denominator \(\sqrt{v_t+\epsilon}\) is tiny.
Then \(\eta/\sqrt{v_t+\epsilon}\) becomes very large \(\Rightarrow\) unstable updates.
We show this: divergence at \(\eta=1\) in our LR sweep.
\end{qa}

% ==================================================
\section{12. Adam}

\begin{qa}{Q: How does Adam work?}
\oral Combines momentum (1st moment) with adaptive scaling (2nd moment) + bias correction.

\textbf{Step 1:} \(m_t = \beta_1 m_{t-1} + (1-\beta_1)g_t\), \quad \(v_t = \beta_2 v_{t-1} + (1-\beta_2)g_t^2\)

\textbf{Step 2:} \(\hat m_t = m_t/(1-\beta_1^t)\), \quad \(\hat v_t = v_t/(1-\beta_2^t)\)

\textbf{Step 3:} \(w_{t+1} = w_t - \eta\hat m_t / (\sqrt{\hat v_t}+\epsilon)\)

\(\hat m_t\) = smoothed direction; \(\hat v_t\) = adaptive scaling.
\end{qa}

\begin{qa}{Q: Why is bias correction needed?}
\oral \(m_0=0\), \(v_0=0\) \(\Rightarrow\) early averages biased toward zero.
\((1-\beta^t)^{-1}\) corrects this.
At \(t=1\): \(1-\beta_2 = 0.001\), so correction factor is \(\times 1000\).
\end{qa}

\begin{qa}{Q: Why can Adam fail to converge?}
\oral Exponential forgetting in \(v_t\) can make \(\hat v_t\) decrease.
If too small: \(\eta/\sqrt{\hat v_t}\) explodes.
\textbf{Fix:} AMSGrad forces \(\hat v_t = \max(\hat v_{t-1}, v_t)\) (non-decreasing).
\end{qa}

\begin{qa}{Q: Why is Adam slower on A9a but faster on MNIST?}
\oral \textbf{A9a (simple convex):} Gradient already correct direction. Momentum (\(\beta_1=0.9\)) dampens it --- overhead. \(\beta_2=0.999\) adapts too slowly.\\
\textbf{MNIST (harder):} 784 dims, 10 classes, noisy gradients. Momentum smooths noise. Per-coordinate scaling matters with heterogeneous features.
\end{qa}

% ==================================================
\section{13. Results Interpretation}

\begin{qa}{Q: Why do adaptive methods converge faster?}
\oral They rescale using gradient history --- avoids too-large or too-small steps per direction.
In our plots: mainly improves early loss decrease.
\end{qa}

\begin{qa}{Q: Why are final accuracies similar?}
\oral Same convex objective \(\Rightarrow\) same global minimum.
Optimizer changes path/speed, not destination.
A9a: all \(\sim\)85\%. MNIST: all \(\sim\)92\%.
\end{qa}

\begin{qa}{Q: What does the gradient norm plot tell us?}
\oral Shows proximity to stationary point.
Adaptive methods bring down large initial gradients faster (normalization).
On A9a: norms don't reach zero --- that's the \textbf{stochastic} gradient norm (irreducible mini-batch variance).
\end{qa}

\begin{qa}{Q: Are your results trustworthy?}
\oral 3 seeds, mean \(\pm\) std shown.
Trends consistent across datasets/losses.
Final report: more seeds + sensitivity + wall-clock.
\end{qa}

% ==================================================
\section{14. Experimental Setup}

\begin{qa}{Q: What model? Neural network?}
\oral \textbf{No.} Linear models.
A9a: logistic regression (binary).
MNIST: softmax regression (10 classes, no hidden layers).
Isolates optimizer effect.
\end{qa}

\begin{qa}{Q: Why not neural networks?}
\oral Introduce confounds (architecture, init, activations).
Goal: clean optimizer comparison.
Linear models = interpretable optimization behavior.
\end{qa}

\begin{qa}{Q: Why A9a and MNIST?}
\oral Two regimes: sparse binary (123 features) vs.\ dense multi-class (784 features).
Tests per-coordinate scaling in different settings.
\end{qa}

\begin{qa}{Q: Did you implement from scratch?}
\oral Yes. All optimizers + losses in NumPy.
Handles dense and sparse matrices (\texttt{scipy.sparse}).
No PyTorch/TensorFlow --- full control.
\end{qa}

% ==================================================
\section{15. Theory vs Practice}

\begin{qa}{Q: Do empirical results match theory?}
\oral Qualitatively yes. Adaptive = faster, as predicted.
But empirical rates much faster than worst-case bounds.
Our objectives have more structure than general theory assumes.
\end{qa}

\begin{qa}{Q: Why care about rates if final accuracy same?}
\oral Compute budget. 500 vs.\ 2000 iterations = \(4\times\) speedup.
For large-scale problems: days vs.\ weeks.
\end{qa}

% ==================================================
\section{16. Outlook}

\begin{qa}{Q: Why LR sweep?}
\oral Shows robustness. Adaptive methods should reduce sensitivity to \(\eta\).
We verify: Adagrad robust, RMSprop fragile.
\end{qa}

\begin{qa}{Q: Why wall-clock comparison?}
\oral Adaptive methods have overhead per step (accumulator updates).
Need to check: still faster in real seconds?
\end{qa}

\begin{qa}{Q: Why add L-BFGS?}
\oral Quasi-Newton method from different family (2nd-order approximation).
Reference beyond 1st-order adaptive methods.
Planned for final report.
\end{qa}

% ==================================================
\section{17. Formula Summary}

\begin{tcolorbox}[title=\textbf{All Updates at a Glance}, colback=unibaselteal!12]
\textbf{SGD:} \(w_{t+1}=w_t-\eta g_t\)

\textbf{Adagrad:} \(w_{t+1} = w_t - \frac{\eta}{\sqrt{G_t}+\epsilon}\, g_t\), \;  \(G_t=G_{t-1}+g_t \odot g_t\)

\textbf{AdaGrad-Norm:} \(w_{t+1} = w_t - \frac{\eta}{b_t+\epsilon}\, g_t\), \; \(b_t^2=b_{t-1}^2+\|g_t\|^2\)

\textbf{RMSprop:} \(w_{t+1} = w_t - \frac{\eta}{\sqrt{v_t}+\epsilon}\, g_t\), \; \(v_t=\rho v_{t-1}+(1-\rho)g_t^2\)

\textbf{Adam:} \(w_{t+1} = w_t - \frac{\eta\hat m_t}{\sqrt{\hat v_t}+\epsilon}\)

\textbf{Polyak:} \(x_{t+1} = x_t - \frac{f(x_t)-f^*}{\|g_t\|^2}\, g_t\) \quad (needs \(f^*\))
\end{tcolorbox}

% ==================================================
\section{18. Quick-Fire Definitions}

\begin{tcolorbox}[title=\textbf{Fast Definitions}, colback=white]
\textbf{Positive definite:} all eigenvalues \(>0\); \(x^\top\!Ax > 0\;\forall x\neq 0\).\\
\textbf{Positive semidefinite:} all eigenvalues \(\geq 0\); \(x^\top\!Ax \geq 0\;\forall x\).\\
\textbf{Lipschitz gradient:} \(\|\nabla f(x)-\nabla f(y)\| \leq L\|x-y\|\).\\
\textbf{Unbiased gradient:} \(\mathbb{E}[g_t] = \nabla f(w_t)\).\\
\textbf{Momentum:} running average of past gradients.\\
\textbf{Bias correction:} \((1-\beta^t)^{-1}\); compensates zero init.\\
\textbf{Per-coordinate:} each feature gets own effective LR.\\
\textbf{Forgetting:} discards old info (RMSprop \(\rho\), Adam \(\beta_2\)).\\
\textbf{Sublinear:} \(O(1/\sqrt{T})\); error goes to zero slowly.\\
\textbf{Linear:} \(e_t \leq \rho^t\); straight line on log-plot.\\
\textbf{Superlinear:} faster than linear; Newton near optimum.\\
\textbf{Strongly convex:} \(\nabla^2 f \succeq \mu I\); unique minimum.
\end{tcolorbox}

% ==================================================
\section{19. Emergency Answers}

\begin{qa}{If you don't know}
\oral We did not investigate this specific point yet, but based on the theory I would expect \dots
\end{qa}

\begin{qa}{If they want more theory}
\oral The formal guarantee depends on smoothness and bounded variance. Intuitively, the method rescales updates based on gradient history.
\end{qa}

\begin{qa}{If they criticize results}
\oral That's a fair point. These are progress results. Final report: LR sensitivity, wall-clock time, stronger \(\lambda\), more seeds.
\end{qa}

\begin{qa}{Are adaptive always better?}
\oral No. Faster early, but can need tuning and have weaker theory. Our results: speed yes, final accuracy no.
\end{qa}

% ==================================================
\section{20. Ten Sentences to Memorize}

\begin{tcolorbox}[title=\textbf{Must-Know}, colback=unibaselteal!12]
\begin{enumerate}
    \item Adaptive methods improve \textbf{speed}, not final accuracy.
    \item We use \textbf{linear models} to isolate optimizer effect.
    \item Adagrad: \textbf{accumulated gradient magnitude}, per-coordinate.
    \item RMSprop: \textbf{forgets old gradients}, adapts quickly.
    \item Adam = momentum + adaptive; theory is \textbf{delicate}.
    \item Final accuracies similar: same objective, same minimum.
    \item Non-convex regularizer: limited at moderate \(\lambda\).
    \item A9a (sparse, 123d) vs.\ MNIST (dense, 784d).
    \item Worst-case bounds are conservative; empirics faster.
    \item Report: LR sweep, wall-clock, stronger \(\lambda\), L-BFGS.
\end{enumerate}
\end{tcolorbox}

\end{multicols}

\newpage

% ==================================================
\section*{Ultra-Short Version for the Table}

\begin{multicols}{2}

\begin{tcolorbox}[title=\textbf{Main Message}, colback=unibaselteal!15]
Adaptive methods improve \textbf{speed}, not necessarily final accuracy.
\end{tcolorbox}

\begin{tcolorbox}[title=\textbf{Why Adaptive?}, colback=white]
Fixed \(\eta\) is hard. Adaptive methods use gradient history to rescale updates.
\end{tcolorbox}

\begin{tcolorbox}[title=\textbf{SGD}, colback=white]
Simple baseline, cheap, but needs careful \(\eta\) and has no coordinate adaptation.
\end{tcolorbox}

\begin{tcolorbox}[title=\textbf{Adagrad}, colback=white]
Per-coordinate scaling via accumulated gradient magnitude. Fast early, but can become conservative.
\end{tcolorbox}

\begin{tcolorbox}[title=\textbf{AdaGrad-Norm}, colback=white]
One global scalar instead of per-coordinate. Simpler, clean non-convex theory.
\end{tcolorbox}

\begin{tcolorbox}[title=\textbf{RMSprop}, colback=white]
Forgets old gradients. Fast and practical, but no formal guarantee. Diverges at high \(\eta\).
\end{tcolorbox}

\begin{tcolorbox}[title=\textbf{Adam}, colback=white]
Momentum + adaptive scaling. Powerful, but convergence more delicate (AMSGrad fix).
\end{tcolorbox}

\begin{tcolorbox}[title=\textbf{Polyak Step Size}, colback=white]
\(\eta_t = (f(x_t)-f^*)/\|g_t\|^2\). Function-value-based. Needs \(f^*\). Not in our experiments.
\end{tcolorbox}

\begin{tcolorbox}[title=\textbf{A9a Result}, colback=white]
RMSprop and Adagrad fastest; all reach \(\sim\)85\% test accuracy.
\end{tcolorbox}

\begin{tcolorbox}[title=\textbf{MNIST Result}, colback=white]
Adagrad fastest; per-coordinate scaling crucial in 784-dim setting. All reach \(\sim\)92\%.
\end{tcolorbox}

\begin{tcolorbox}[title=\textbf{Non-Convex Regularizer}, colback=white]
\(\lambda=0.01\): mild non-convexity, little effect. \(\lambda \geq 0.1\): underfitting.
\end{tcolorbox}

\begin{tcolorbox}[title=\textbf{Why Linear Models?}, colback=white]
Controlled setup. Isolates optimizer. Convex = clean theory.
\end{tcolorbox}

\begin{tcolorbox}[title=\textbf{Logistic Loss (A9a)}, colback=white]
\(f(w) = \frac{1}{n}\sum_i \log(1+\exp(-y_i w^\top x_i))\). Convex. Binary.
\end{tcolorbox}

\begin{tcolorbox}[title=\textbf{Softmax Loss (MNIST)}, colback=white]
\(f(W) = -\frac{1}{n}\sum_i \log \frac{\exp(w_{y_i}^\top x_i)}{\sum_k \exp(w_k^\top x_i)}\). Convex. 10 classes.
\end{tcolorbox}

\begin{tcolorbox}[title=\textbf{Outlook}, colback=white]
LR sweep, wall-clock time, stronger \(\lambda\), L-BFGS baseline.
\end{tcolorbox}

\end{multicols}

\end{document}
